\documentclass[10pt,a4paper]{article}
\usepackage{amsmath,amssymb,amsthm,graphicx}
% Minimal editorial layout adapter for the pinned JEP source
% sources/jep-2022-195-contact-systolic/pinned-original-native.tex.
% The publisher document class (cedram) is not part of the retained material, and it
% is the class that used to define a handful of plain text abbreviations used by the
% author's prose. This adapter supplies exactly those abbreviations and nothing else:
% no mathematics, no notation, no theorem environment, no counter.
% HoTT/JEP retention convention: CC BY 4.0, attribution to the original authors.
\providecommand{\ie}{i.e.\ }
\providecommand{\eg}{e.g.\ }
\providecommand{\skpt}{\smallskip}
\usepackage[scr=rsfs,cal=euler]{mathalfa}
\let\mathds\mathbb
\let\epsilon\varepsilon
\newcommand{\Changel}{\mathcode`l="7160}
\newcommand{\Changelback}{\mathcode`l="716C}
\newcommand{\NoChangel}[1]{%
\expandafter\let\csname old\string#1\endcsname=#1
\let#1=\relax
\newcommand{#1}{\mathcode`l="716C\csname old\string#1\endcsname\mathcode`l="7160 }%
}

\NoChangel{\log}
\NoChangel{\ln}
\NoChangel{\lim}
\NoChangel{\limsup}
\NoChangel{\liminf}
\NoChangel{\varlimsup}
\NoChangel{\varliminf}
\NoChangel{\varinjlim}
\NoChangel{\varprojlim}

\newenvironment{enumeratei}
{\bgroup\def\theenumi{\roman{enumi}}\def\theenumii{\arabic{enumii}}\begin{enumerate}}
{\end{enumerate}\egroup}
\newcommand{\Psfrac}[2]{(\sfrac{#1}{#2})}
\newcommand{\psfrac}[2]{\sfrac{(#1)}{#2}}
\newcommand{\spfrac}[2]{\sfrac{#1}{(#2)}}
\newcommand\mto{\mathchoice{\longmapsto}{\mapsto}{\mapsto}{\mapsto}}
\newcommand\hto{\mathrel{\lhook\joinrel\to}}

\usepackage{graphicx}

\newtheorem{thm}{Theorem}[section]
\newtheorem{cor}[thm]{Corollary}
\newtheorem{lem}[thm]{Lemma}
\newtheorem{prop}[thm]{Proposition}
\newtheorem{quest}{Question}[section]
\newtheorem{conj}{Conjecture}[section]
\newtheorem{maintheorem}{Theorem}
\theoremstyle{definition}
\newtheorem{rem}[thm]{Remark}
\newtheorem{exm}[thm]{Example}
\newtheorem{ass}[thm]{Assumption}
\renewcommand*{\themaintheorem}{\Alph{maintheorem}}

\DeclareMathOperator{\area}{area}
\DeclareMathOperator{\lcm}{lcm}
\DeclareMathOperator{\const}{const}
\DeclareMathOperator{\Cal}{Cal}
\DeclareMathOperator{\Flux}{\mathrm{Flux}}
\newcommand{\id}{\mathrm{id}}
\newcommand{\Ham}{\mathrm{Ham}}
\newcommand{\std}{\mathrm{std}}
\newcommand{\vol}{\mathrm{vol}}
\newcommand{\fix}{\mathrm{fix}}
\newcommand{\pr}{\mathrm{pr}}
\newcommand{\Tan}{\mathrm{T}}
\newcommand{\N}{\mathds{N}}
\newcommand{\Z}{\mathds{Z}}
\newcommand{\R}{\mathds{R}}
\newcommand{\CC}{\mathds{C}}\let\C\CC
\newcommand{\UU}{\mathcal{U}}
\newcommand{\VV}{\mathcal{V}}

\let\eul e
\let\diff d
\title{Weinstein tubular neighborhood for the contact form, with an explicit collar model}
\author{The Univalent Foundations Program}
\date{doi:10.5802/jep.195; CC BY-SA 3.0 Unported}
\begin{document}
\maketitle
\setcounter{section}{2}
\setcounter{thm}{9}
\noindent\textit{Editorial note.} The counted statement is the author's labelled lemma `l:Weinstein` (source lines 715-738); the proof environment starts after one blank line, so the pairing is by adjacency.\par\medskip
\noindent\textit{Editorial note.} The retained context is the author's own displayed equation for the normalised primitive on the collar, which the statement itself invokes by label.\par\medskip
\noindent\textit{Editorial note.} The retained counter values (section 2, theorem offset 9) make the counted lemma print as Lemma 2.10; this article carries no label-number dump, so the number was inferred by counting the theorem-class environments of section 2 and is reported with that caveat.\par\medskip
\noindent\textit{Editorial note.} No retained line refers to a label, citation or figure outside the retained spans, so the package declares no external reference.\par\medskip
\noindent\textit{Editorial note.} The retained preamble is a minimal editorial adapter supplying only the plain text abbreviations that the publisher class used to define, followed by the native author macro block of the pinned source.\par\medskip
\noindent\textit{Editorial scope.} One original result of the pinned publisher article, counted once: the Weinstein tubular neighbourhood theorem for the contact form with its collar model (author label l:Weinstein).\par\medskip
\noindent\textit{Allowed standard prerequisites (editorial).}\begin{itemize}
\item Canonical Liouville form on the cotangent bundle. lambda\_std is the standard Liouville one-form on T*Sigma, with d lambda\_std the canonical symplectic form.. It is the target form whose pull-back is computed at source lines 754-761 and 802.
\item Collar coordinates and the normalised primitive. On the collar A\_C = [0,rho) x S\textasciicircum{}1 the primitive nu\_0 has the displayed form with a constant a\_C.. It is used to write f\_0 explicitly at source lines 765-772.
\item Moser's trick. Given a family of forms mu\_t with mu\_0 = mu\_1 up to exactness one solves X\_t interior dmu\_t = -(d/dt)mu\_t and integrates to an isotopy pulling mu\_1 back to mu\_0 up to an exact form.. It is the whole second half of the proof, source lines 777-808.
\end{itemize}\medskip
\begin{equation}
\label{e:lambda|_A}
\nu_0|_{A_C} = \bigl( a_C - \tfrac{1}{2} r^2 \bigr)\, ds,
\end{equation}

\begin{lem}[Weinstein tubular neighborhood]
\label{l:Weinstein}
There exists an open neighborhood $U\subset\Sigma\times\Sigma$ of the diagonal $\Delta_\Sigma=\{(z,z)\mid z\in\Sigma\}$, an open neighborhood \hbox{$V\subset\Tan^*\Sigma$} of the $0$-section $0_\Sigma\subset\Tan^*\Sigma$, and a smooth map $\psi:U\to V$ that restricts to a diffeomorphism $\psi:U\cap\mathrm{int}(\Sigma\times\Sigma)\to V\cap \Tan^* \mathrm{int} (\Sigma)$, and satisfies
\[
\psi(\Delta_\Sigma)=0_\Sigma,
\qquad
\psi^*\lambda_{\std}=\nu_2-\nu_1 + \diff f,
\] where $f:\Sigma\times\Sigma\to\R$ is a smooth function such that $f|_{\Delta_\Sigma}\equiv0$ and
\[
\diff f(z_1,z_2)=(\nu_1-\nu_2)_{(z_1,z_2)},
\qquad
f(z_1,z_2)=-\int_{[z_1,z_2]} \!\!\!\!\!\!\nu_0,
\]
for every pair $(z_1,z_2)\in U\cap(\partial\Sigma\times\partial\Sigma)$.
If $A_C\equiv[0,\rho)\times S^1$ is the collar neighborhood of a connected component $C$ of $\partial\Sigma$ on which $\nu_0$ has the form \eqref{e:lambda|_A}, the restriction $\psi|_{U\cap(A_C\times A_C)}$ has the form
\begin{equation}
\label{psi}
\begin{split}
\psi(r,s,R,S)=\big( R,s,\,R\,(S-s),\tfrac 1 2 (r^2-R^2) \big),
\\
\forall (r,s,R,S)\in U\cap(A_C\times A_C).
\end{split}
\end{equation}
\end{lem}

\begin{proof}
We first provide the construction within the collar neighborhood $A=A_C$ of each connected component $C$ of $\partial\Sigma$.
We consider a small enough neighborhood $W\subset A\times A$ of the diagonal $\Delta_A=\{(z,z)\mid z\in A\}$ so that, for each $(r,s,R,S)\in W$, the points $s,S\in S^1$ are not antipodal. We define the map
\[
\kappa_0:W\to \underbrace{\Tan^*A}_{A\times \R^2},
\qquad
\kappa_0(r,s,R,S)=\big( R,s,\,R\,(S-s),\tfrac 12 (r^2-R^2) \big).
\]
This map restricts to a diffeomorphism onto its image
\[
\kappa_0:W\cap\mathrm{int}(A\times A)\to \mathrm{int}(\Tan^*A),
\]
and satisfies $\kappa_0(\Delta_A)=0_A$ and
\begin{align*}
\kappa_0^*\lambda_{\std}
&=
R\,(S-s)\,\diff R + \tfrac 12(r^2-R^2)\,\diff s\\
&=
-\tfrac 12\,R^2\,\diff S + \tfrac 12\,r^2\,\diff s + \diff\big( \tfrac 12(S-s)\,R^2 \big)\\
&=
-\nu_1+\nu_2+\diff\big( (a_C -\tfrac 12 R^2)(s-S) \big)\\
&=-\nu_1+\nu_2+\diff f_0,
\end{align*}
where
\[
f_0:W\to\R,\qquad f_0(r,s,R,S)=(a_C -\tfrac 12 R^2)(s-S).
\]
Notice that $f_0|_{\Delta_A}\equiv0$ and
\begin{align*}
f_0(0,s,0,S) & =a_C (s-S) = - \int_{[z_1,z_2]} \nu_0,
\\
\diff f_0(0,s,0,S) & = a_C \,\diff s-a_C\,\diff S = \nu_1-\nu_2,
\end{align*}
where $z_1=(0,s)$ and $z_2=(0,S)$.

For some sufficiently small neighborhood $U\subset\Sigma\times\Sigma$ of the diagonal $\Delta_\Sigma$, we choose an arbitrary smooth function $f_1:U\to\R$ that coincides with $f_0$ on a neighborhood of $U\cap(\partial\Sigma\times\partial\Sigma)$ and satisfies $f_1|_{\Delta_\Sigma}\equiv0$ and $\diff f_1=\nu_1-\nu_2$ at all points of the diagonal $\Delta_\Sigma$. Up to further shrinking the neighborhood $U$, we also choose a smooth map $\kappa_1:U\to\Tan^*\Sigma$ that coincides with $\kappa_0$ on a neighborhood of $W\cap(\partial\Sigma\times\partial\Sigma)$, restricts to a diffeomorphism onto its image $\kappa_1:U\cap\mathrm{int}(\Sigma\times\Sigma)\to\Tan^*(\mathrm{int}(\Sigma))$, and such that $\kappa_1(\Delta_\Sigma)=0_\Sigma$.

We now conclude the proof by means of a typical application of Moser's trick. We~set
\[
\mu_t:=t\kappa_1^*(\lambda_{\std}) + (1-t)(\nu_2-\nu_1),
\]
and we look for an isotopy $\phi_t:U\to\Sigma\times\Sigma$, defined after possibly further shrinking the neighborhood $U$, such that $\phi_0=\id$ and $\phi_t^*\mu_t-\mu_0$ is exact. We denote by $X_t$ the time-dependent vector field generating $\phi_t$ and compute
\begin{align*}
\tfrac{\diff}{\diff t}\phi_t^*\mu_t
&=\phi_t^* ( L_{X_t} \mu_t + \tfrac{\diff}{\diff t} \mu_t) = \phi_t^*\big( X_t\,\lrcorner\,\diff\mu_t + \diff(\mu_t(X_t))+ \mu_1-\mu_0 \big)\\
&=\phi_t^*\big( X_t\,\lrcorner\,\diff\mu_t + \kappa_1^*\lambda_{\std}-\nu_2+\nu_1-\diff f_1 + \diff(\mu_t(X_t)+f_1)\big).
\end{align*}
Notice that the symplectic forms $\kappa_1^*\diff\lambda_{\std}$ and $-\diff\nu_1+\diff\nu_2$ define the same orientation, since they coincide on a neighborhood of $U\cap(\partial\Sigma\times\partial\Sigma)$. Therefore $\diff\mu_t$ is symplectic away from $\partial(\Sigma\times \Sigma)$ for every $t\in[0,1]$. We choose the vector field $X_t$ so that
\[
X_t\,\lrcorner\,\diff\mu_t + \kappa_1^*\lambda_{\std}-\nu_2+\nu_1-\diff f_1
=
0.
\]
Notice that $X_t$ vanishes on the diagonal $\Delta_\Sigma$ and on a neighborhood of $U\cap(\partial\Sigma\times\partial\Sigma)$.
Moreover
\[
\tfrac{\diff}{\diff t}\phi_t^*\mu_t
=
\phi_t^*\diff(\mu_t(X_t)+f_1).
\]
Up to shrinking the neighborhood $U$, we obtain a well defined isotopy $\phi_t:U\to\Tan^*\Sigma$ that coincides with $\kappa_0$ on a neighborhood of $U\cap(\partial\Sigma\times\partial\Sigma)$ and satisfies
\[
\phi_1^*\kappa_1^*(\lambda_{\std})=-\nu_1+\nu_2 + \diff f,
\]
where
\[
f(w)=\int_0^1 \big(\mu_t(X_t)+f_1\big)\circ\phi_t(w)\,\diff t.
\]
The desired map is $\psi:=\kappa_1\circ\phi_1$.
\end{proof}

\end{document}
